% ============================================================
\chapter*{Pr\'eface}
\addcontentsline{toc}{chapter}{Pr\'eface}
\markboth{Pr\'eface}{Pr\'eface}
% ============================================================

L'analyse topologique des donn\'ees (ATD, ou \emph{Topological Data Analysis} --- TDA
en anglais) est un domaine relativement jeune, n\'e au d\'ebut des ann\'ees~2000
\`a la confluence de la topologie alg\'ebrique, de la g\'eom\'etrie computationnelle
et de la science des donn\'ees. Son objectif fondamental est d'extraire des
informations \emph{qualitatives} et \emph{structurelles} \`a partir de jeux de
donn\'ees complexes, en exploitant des invariants topologiques robustes au bruit
et aux d\'eformations continues.

\medskip

Alors que les m\'ethodes statistiques classiques reposent sur des mod\`eles
param\'etriques ou sur des hypoth\`eses de distribution, la TDA adopte une
perspective radicalement diff\'erente~: elle cherche \`a caract\'eriser la
\emph{forme} des donn\'ees. Cette approche s'est r\'ev\'el\'ee particuli\`erement
fructueuse dans des domaines aussi vari\'es que la biologie mol\'eculaire,
les neurosciences, la science des mat\'eriaux, la finance quantitative et
l'apprentissage profond.

%% ----------------------------------------------------------------
\section*{Objectifs de ce cours}

Ce cours est con\c{c}u pour des \'etudiants de niveau Master~2 ou doctorat,
disposant d'une formation de base en~:
\begin{itemize}
  \item \textbf{Alg\`ebre et topologie}~: groupes, anneaux, modules~;
        espaces topologiques, compacit\'e, connexit\'e.
  \item \textbf{Analyse et probabilit\'es}~: espaces m\'etriques, th\'eorie de
        la mesure, convergence.
  \item \textbf{Algorithmique et programmation}~: complexit\'e, structures de
        donn\'ees, Python.
\end{itemize}

\`A l'issue de ce cours, l'\'etudiant sera capable~:
\begin{enumerate}
  \item De construire et manipuler les principaux complexes simpliciaux
        (Vietoris--Rips, \v{C}ech, Alpha) associ\'es \`a un nuage de points.
  \item De calculer l'homologie simpliciale et l'homologie persistante,
        tant sur le plan th\'eorique qu'algorithmique.
  \item D'interpr\'eter les diagrammes de persistance et les codes-barres,
        et d'en \'evaluer la stabilit\'e.
  \item D'appliquer l'algorithme Mapper pour la visualisation exploratoire.
  \item D'int\'egrer les descripteurs topologiques dans des pipelines
        d'apprentissage automatique et d'apprentissage profond.
  \item De mener un projet de recherche appliquant la TDA \`a des donn\'ees
        r\'eelles.
\end{enumerate}

%% ----------------------------------------------------------------
\section*{Organisation du cours}

Le cours est structur\'e en onze chapitres, regroup\'es en quatre parties~:

\begin{description}
  \item[Partie~I --- Fondements (Chapitres~1--3)]
    Motivations, complexes simpliciaux, homologie et introduction
    \`a l'homologie persistante.
  \item[Partie~II --- Th\'eorie de la persistance (Chapitres~4--6)]
    Diagrammes de persistance, th\'eor\`emes de stabilit\'e et construction
    des complexes filtr\'es.
  \item[Partie~III --- M\'ethodes et algorithmes (Chapitres~7--8)]
    Algorithme Mapper, distances entre diagrammes de persistance
    (Wasserstein, bottleneck).
  \item[Partie~IV --- TDA et apprentissage (Chapitres~9--11)]
    Int\'egration avec le Machine Learning, applications concr\`etes
    et liens avec le Deep Learning.
\end{description}

Chaque chapitre contient des d\'efinitions rigoureuses, des th\'eor\`emes avec
preuves (ou esquisses de preuves), des exemples d\'etaill\'es, des
impl\'ementations en Python utilisant \texttt{gudhi}, \texttt{ripser} et
\texttt{giotto-tda}, ainsi que des exercices de difficult\'e croissante.

%% ----------------------------------------------------------------
\section*{Notations et conventions}

\begin{center}
\renewcommand{\arraystretch}{1.3}
\begin{tabular}{cl}
\toprule
\textbf{Symbole} & \textbf{Signification} \\
\midrule
$\R^d$ & Espace euclidien de dimension $d$ \\
$\N, \Z, \Q$ & Entiers naturels, entiers relatifs, rationnels \\
$\mathbb{F}$ & Corps fini (souvent $\mathbb{F}_2 = \Z/2\Z$) \\
$K$ & Complexe simplicial abstrait \\
$\sigma = [v_0, \ldots, v_p]$ & $p$-simplexe \\
$C_p(K)$ & Groupe des $p$-cha\^{\i}nes de $K$ \\
$\partial_p$ & Op\'erateur bord en dimension $p$ \\
$H_p(K)$ & $p$-i\`eme groupe d'homologie de $K$ \\
$\beta_p$ & $p$-i\`eme nombre de Betti ($= \dim H_p$) \\
$\mathrm{VR}(X, r)$ & Complexe de Vietoris--Rips au param\`etre $r$ \\
$\check{C}(X, r)$ & Complexe de \v{C}ech au param\`etre $r$ \\
$\mathrm{Dgm}_p$ & Diagramme de persistance en dimension $p$ \\
$d_B, W_p$ & Distance bottleneck, distance de Wasserstein d'ordre $p$ \\
$\norm{\cdot}$ & Norme (euclidienne par d\'efaut) \\
$\abs{S}$ & Cardinal de l'ensemble $S$ \\
$\inner{\cdot}{\cdot}$ & Produit scalaire \\
\bottomrule
\end{tabular}
\end{center}

%% ----------------------------------------------------------------
\section*{Environnement logiciel}

Les impl\'ementations propos\'ees dans ce cours utilisent
\textbf{Python~3.10+} avec les biblioth\`eques suivantes~:

\begin{itemize}
  \item \textbf{GUDHI} (\url{https://gudhi.inria.fr})~: biblioth\`eque
        C++/Python d\'evelopp\'ee par l'INRIA, offrant un large \'eventail de
        structures simpliciales, d'algorithmes de persistance et de
        fonctionnalit\'es statistiques.
  \item \textbf{Ripser} (\url{https://ripser.scikit-tda.org})~:
        impl\'ementation ultra-rapide de l'homologie persistante pour les
        complexes de Vietoris--Rips.
  \item \textbf{giotto-tda} (\url{https://giotto-ai.github.io/gtda-docs})~:
        biblioth\`eque orient\'ee Machine Learning, compatible scikit-learn.
  \item \textbf{scikit-learn}, \textbf{NumPy}, \textbf{SciPy},
        \textbf{Matplotlib}~: pr\'etraitement, calcul num\'erique et
        visualisation.
  \item \textbf{PyTorch}~: pour les mod\`eles d'apprentissage profond
        int\'egrant des couches topologiques.
\end{itemize}

\begin{attention}[title=\textbf{Installation recommand\'ee}]
Nous recommandons l'utilisation d'un environnement virtuel~:
\begin{verbatim}
python -m venv tda-env
source tda-env/bin/activate
pip install gudhi ripser giotto-tda scikit-learn torch matplotlib
\end{verbatim}
\end{attention}

%% ----------------------------------------------------------------
\section*{Rep\`eres historiques}

L'analyse topologique des donn\'ees puise ses racines dans des travaux
fondateurs~:

\begin{itemize}
  \item \textbf{Ann\'ees 1990}~: Edelsbrunner, Letscher et Zomorodian
        introduisent les premi\`eres notions d'homologie persistante.
  \item \textbf{2002}~: Zomorodian et Carlsson formalisent l'homologie
        persistante et son algorithme de calcul via la d\'ecomposition en
        modules gradu\'es.
  \item \textbf{2005}~: Cohen-Steiner, Edelsbrunner et Harer d\'emontrent
        le th\'eor\`eme de stabilit\'e fondamental pour les diagrammes de
        persistance.
  \item \textbf{2007}~: Singh, M\'emoli et Carlsson introduisent
        l'algorithme Mapper.
  \item \textbf{2009}~: Chazal, Cohen-Steiner, Glisse, Guibas et Oudot
        \'etendent les r\'esultats de stabilit\'e aux distances de
        Wasserstein.
  \item \textbf{2015--pr\'esent}~: Int\'egration croissante de la TDA
        avec le Machine Learning et le Deep Learning.
\end{itemize}

%% ----------------------------------------------------------------
\section*{Remerciements}

Ce cours doit beaucoup aux travaux pionniers de Gunnar Carlsson, Herbert
Edelsbrunner, Fr\'ed\'eric Chazal, Steve Oudot, et de toute la communaut\'e
TDA. Nous remercions \'egalement les d\'eveloppeurs des biblioth\`eques
GUDHI, Ripser et giotto-tda pour la qualit\'e de leurs outils.

\begin{intuition}
La TDA repose sur une id\'ee simple mais profonde~: \emph{la forme des
donn\'ees contient de l'information}. Les outils de la topologie
alg\'ebrique permettent de quantifier cette forme de mani\`ere rigoureuse,
stable et calculable. Ce cours vous donnera les cl\'es pour exploiter
cette id\'ee dans vos propres recherches.
\end{intuition}

\bigskip
\hfill\textit{Les auteurs, mars 2026}
